% ===== Preâmbulo compartilhado (Tectonic / XeTeX) — documentos da FTL094 =====
\documentclass[11pt,a4paper]{article}
\usepackage{fontspec}
\usepackage[brazil]{babel}
\usepackage{geometry}
\geometry{a4paper,margin=2.5cm}
\usepackage{amsmath,amssymb}
\usepackage{graphicx}
\usepackage{booktabs}
\usepackage{array}
\usepackage{longtable}
\usepackage{tabularx}
\usepackage{float}
\usepackage{enumitem}
\usepackage{xcolor}
\usepackage{tikz}
\usetikzlibrary{arrows.meta,positioning,automata,shapes.geometric,fit,backgrounds,calc}
\usepackage{pgfplots}
\pgfplotsset{compat=1.18}
\usepackage{listings}
\usepackage{caption}
\usepackage{fancyhdr}
\usepackage[hidelinks]{hyperref}

% ----- cores -----
\definecolor{ufamblue}{RGB}{0,51,102}
\definecolor{codegreen}{rgb}{0,0.5,0}
\definecolor{codeblue}{rgb}{0,0,0.7}
\definecolor{codered}{rgb}{0.6,0,0}
\definecolor{lgray}{rgb}{0.95,0.95,0.95}

% ----- listagens C -----
\lstdefinestyle{c}{language=C,basicstyle=\ttfamily\footnotesize,
  keywordstyle=\color{codeblue}\bfseries,commentstyle=\color{codegreen}\itshape,
  stringstyle=\color{codered},numbers=left,numberstyle=\tiny\color{gray},
  frame=single,breaklines=true,tabsize=2,showstringspaces=false,columns=fullflexible}

% ----- constantes do projeto -----
\newcommand{\projnome}{Sistema de Navegação Autônoma para o Robô Móvel Khepera~IV}
\newcommand{\disciplina}{FTL094 --- Engenharia de Software de Tempo Real}
\newcommand{\periodo}{Período Letivo 2026/1}
\newcommand{\versaodoc}{1.0}

% ----- cabeçalho/rodapé -----
\pagestyle{fancy}\fancyhf{}
\rhead{\small \tipodoc}
\lhead{\small Khepera~IV --- FTL094}
\cfoot{\thepage}
\renewcommand{\headrulewidth}{0.4pt}

% ----- coluna X centralizada p/ tabularx -----
\newcolumntype{C}{>{\centering\arraybackslash}X}

% ----- capa padronizada -----
\newcommand{\fazcapa}{%
\begin{titlepage}\centering
{\large\bfseries UNIVERSIDADE FEDERAL DO AMAZONAS (UFAM)}\\[3pt]
{\normalsize Centro de P\&D em Tecnologia Eletrônica e da Informação --- CETELI}\\[2pt]
{\normalsize Programa de Pós-Graduação em Engenharia Elétrica --- PPGEE}\\[26pt]
{\large \disciplina}\\[2pt]
{\normalsize \periodo}\\[64pt]
\rule{\textwidth}{1.2pt}\\[12pt]
{\Huge\bfseries\color{ufamblue} \tipodoc}\\[14pt]
{\Large \projnome}\\[10pt]
\rule{\textwidth}{1.2pt}\\[48pt]
{\large\bfseries Equipe}\\[8pt]
{\normalsize Gabriel Henrique de Souza Nazaré}\\[3pt]
{\normalsize Luan Alexandre Ramos Barreto}\\[3pt]
{\normalsize Matheus Rocha Canto}\\[3pt]
{\normalsize Matheus Serrão Uchôa}\\[3pt]
{\normalsize Ricardo Mendonça Braz}\\[3pt]
{\normalsize Yan Fernandes da Silva}\\[40pt]
{\normalsize Docente responsável: Prof. Dr.~[nome do docente]}\\
\vfill
{\small Documento de tipo \emph{\tipodoc} --- Versão \versaodoc}\\[2pt]
{\small Manaus --- AM, \datadoc}
\end{titlepage}}

% ----- tabela de controle de versões -----
\newcommand{\controleversoes}{%
\section*{Controle de versões}
\begin{tabularx}{\textwidth}{@{}l l l X@{}}
\toprule
\textbf{Versão} & \textbf{Data} & \textbf{Autor} & \textbf{Descrição} \\
\midrule
1.0 & \datadoc & Equipe FTL094 & Emissão inicial do documento. \\
\bottomrule
\end{tabularx}
\bigskip}

\newcommand{\tipodoc}{Plano de Projeto}
\newcommand{\datadoc}{16 de junho de 2026}

\begin{document}
\fazcapa
\controleversoes
\tableofcontents
\newpage

% =====================================================================
\section{Introdução}
% =====================================================================
\subsection{Objetivo do documento}
Este Plano de Projeto define a organização, o processo de desenvolvimento, o cronograma, os recursos e a gerência de riscos do projeto \emph{\projnome}, desenvolvido como trabalho final da disciplina \disciplina. O documento serve de instrumento de planejamento e controle, orientando a equipe e estabelecendo o compromisso com o cliente (a coordenação da disciplina).

\subsection{Escopo do projeto}
O projeto consiste no desenvolvimento do \textbf{software de controle de um robô móvel autônomo} capaz de deslocar-se de uma posição inicial a uma posição-alvo, em um ambiente povoado por obstáculos, \emph{sem colisões}, em condições representativas de competições de robótica. A plataforma-alvo é o robô de tração diferencial \textbf{Khepera~IV} (acervo do CETELI/UFAM); o desenvolvimento e a validação são realizados no simulador \textbf{Webots~R2025a}, garantindo reprodutibilidade antes de eventual transferência para o robô físico. Estão \emph{fora do escopo} a construção de hardware, a navegação em ambientes dinâmicos com obstáculos móveis e a construção de mapas globais (SLAM).

\subsection{Definições, acrônimos e abreviações}
\begin{description}[style=nextline,leftmargin=1.2cm,font=\ttfamily]
  \item[FSM] Máquina de estados finitos (\emph{Finite State Machine}).
  \item[IR] Sensor de proximidade infravermelho.
  \item[IMU] Unidade de medição inercial.
  \item[RT] Tempo real (\emph{Real-Time}).
  \item[WCET] Tempo de execução de pior caso.
  \item[EAP] Estrutura analítica do projeto.
\end{description}

\subsection{Referências}
\begin{enumerate}[label={[}\arabic*{]},leftmargin=1.2cm,noitemsep]
  \item I.~Sommerville, \emph{Engenharia de Software}, 10ª ed. Pearson, 2019.
  \item R.~S.~Pressman, \emph{Engenharia de Software: Uma Abordagem Profissional}, 8ª ed. AMGH, 2016.
  \item Documento de Análise de Requisitos do projeto (artefato companheiro).
  \item R.~Siegwart \emph{et al.}, \emph{Introduction to Autonomous Mobile Robots}, MIT Press, 2011.
\end{enumerate}

% =====================================================================
\section{Visão geral do projeto}
% =====================================================================
\subsection{Contexto e justificativa}
A navegação autônoma é problema central da robótica móvel e exige a integração de percepção, decisão e controle sob restrições de tempo real. O presente projeto materializa esses conceitos em um produto de software embarcado, servindo simultaneamente como exercício completo do ciclo de Engenharia de Software (planejamento, requisitos, projeto, implementação, verificação) e como solução técnica a um problema concreto de robótica.

\subsection{Objetivos}
\textbf{Objetivo geral:} entregar um software que conduza o Khepera~IV de um ponto~A a um ponto~B desviando de todos os obstáculos do percurso.

\noindent\textbf{Objetivos específicos (SMART):}
\begin{itemize}[noitemsep,topsep=2pt]
  \item alcançar o alvo (tolerância de 7~cm) em cenário com 5 obstáculos, sem colisões;
  \item garantir laço de controle periódico de 16~ms (restrição de tempo real);
  \item prover robustez contra aprisionamento em mínimos locais;
  \item documentar todos os artefatos exigidos pela disciplina.
\end{itemize}

\subsection{Entregáveis}
\begin{table}[H]\centering
\caption{Entregáveis do projeto e prazos.}
\begin{tabularx}{\textwidth}{@{}l X l@{}}
\toprule
\textbf{Id} & \textbf{Entregável} & \textbf{Prazo} \\
\midrule
E1 & Plano de Projeto + Análise de Requisitos & 16/06/2026 \\
E2 & Documento de Arquitetura/Projeto do Sistema & 30/06/2026 \\
E3 & Plano de Testes e Protocolo de Testes & 30/06/2026 \\
E4 & Código-fonte e entrega final ao cliente & 09/07/2026 \\
E5 & Apresentação final do projeto & 15/07/2026 \\
\bottomrule
\end{tabularx}
\end{table}

\subsection{Premissas e restrições}
\textbf{Premissas:} disponibilidade do simulador Webots; ambiente de testes estático e conhecido; alvo fornecido por coordenadas. \textbf{Restrições:} linguagem \textbf{C} (cadeia de ferramentas embarcada no Webots); sensoriamento de proximidade de curto alcance ($\approx$ 2--20~cm); ausência de mapa global \emph{a priori}.

% =====================================================================
\section{Organização do projeto}
% =====================================================================
\subsection{Equipe e papéis}
\begin{table}[H]\centering
\caption{Estrutura da equipe e responsabilidades.}
\begin{tabularx}{\textwidth}{@{}l l X@{}}
\toprule
\textbf{Integrante} & \textbf{Papel} & \textbf{Responsabilidades} \\
\midrule
Matheus Serrão Uchôa & Líder técnico & Arquitetura, implementação do controlador, integração. \\
Gabriel Henrique de Souza Nazaré & Eng.\ de requisitos & Levantamento e especificação de requisitos. \\
Matheus Rocha Canto & Arquiteto/Projetista & Projeto da arquitetura e decomposição modular. \\
Luan Alexandre Ramos Barreto & Eng.\ de testes & Plano de testes, execução e protocolo. \\
Ricardo Mendonça Braz & Gerente de projeto & Cronograma, riscos e monitoramento. \\
Yan Fernandes da Silva & Documentação/Integração & Documentação técnica e integração dos artefatos. \\
\bottomrule
\end{tabularx}
\end{table}

\subsection{Partes interessadas (\emph{stakeholders})}
Cliente/coordenação da disciplina (define a encomenda e avalia); equipe de desenvolvimento; usuário final hipotético (operador do robô em competição); CETELI/UFAM (provê plataforma).

% =====================================================================
\section{Processo de desenvolvimento}
% =====================================================================
Adota-se um \textbf{modelo de processo incremental e iterativo}, justificado pela natureza experimental do controlador, que demanda ciclos curtos de \emph{implementar--medir--ajustar} (típicos de software de controle e de tempo real). Cada incremento agrega um comportamento (atração ao alvo; desvio reativo; contorno comprometido) e é validado por telemetria antes do próximo. As fases mapeiam-se às unidades da disciplina (requisitos, projeto, implementação, verificação), com revisões ao final de cada entrega.

% =====================================================================
\section{Cronograma e marcos}
% =====================================================================
\subsection{Estrutura analítica do projeto (EAP)}
\begin{enumerate}[leftmargin=1.4cm,noitemsep]
  \item[1.] Gerência \quad (1.1 Plano; 1.2 Riscos; 1.3 Acompanhamento)
  \item[2.] Requisitos \quad (2.1 Levantamento; 2.2 Especificação; 2.3 Rastreabilidade)
  \item[3.] Projeto \quad (3.1 Arquitetura; 3.2 Decomposição; 3.3 Projeto detalhado)
  \item[4.] Implementação \quad (4.1 Mundo Webots; 4.2 Controlador C; 4.3 Telemetria)
  \item[5.] Verificação \quad (5.1 Plano de testes; 5.2 Execução; 5.3 Protocolo)
  \item[6.] Entrega \quad (6.1 Código; 6.2 Apresentação)
\end{enumerate}

\subsection{Marcos (\emph{milestones})}
\begin{figure}[H]\centering
\begin{tikzpicture}[x=1.45cm,y=1cm]
  \draw[ufamblue,very thick,-{Stealth}] (0,0) -- (8.6,0);
  \foreach \x/\d/\l in {0.6/16-06/E1, 3/30-06/{E2,E3}, 6/09-07/E4, 8/15-07/E5}{
    \filldraw[ufamblue] (\x,0) circle(2.4pt);
    \node[above,font=\scriptsize,align=center] at (\x,0.15) {\textbf{\l}};
    \node[below,font=\scriptsize] at (\x,-0.15) {\d};
  }
\end{tikzpicture}
\caption{Linha do tempo dos marcos de entrega.}
\end{figure}

\begin{table}[H]\centering
\caption{Cronograma sintético por fase.}
\begin{tabularx}{\textwidth}{@{}X l l l@{}}
\toprule
\textbf{Fase} & \textbf{Início} & \textbf{Fim} & \textbf{Entregável} \\
\midrule
Requisitos e Plano & 02/06 & 16/06 & E1 \\
Arquitetura/Projeto & 16/06 & 30/06 & E2 \\
Plano de Testes & 16/06 & 30/06 & E3 \\
Implementação & 16/06 & 07/07 & --- \\
Verificação/Validação & 30/06 & 08/07 & (protocolo) \\
Entrega de código & 07/07 & 09/07 & E4 \\
Apresentação final & 09/07 & 15/07 & E5 \\
\bottomrule
\end{tabularx}
\end{table}

% =====================================================================
\section{Estimativas e recursos}
% =====================================================================
\subsection{Recursos}
\begin{table}[H]\centering
\caption{Recursos de hardware e software.}
\begin{tabularx}{\textwidth}{@{}l X@{}}
\toprule
\textbf{Tipo} & \textbf{Recurso} \\
\midrule
Software & Webots R2025a; compilador MinGW-w64 (embarcado); editor; Git; \LaTeX{} (Tectonic). \\
Hardware (sim.) & Modelo Khepera~IV (8 IR, 5 ultrassom, câmera, IMU, tração diferencial). \\
Hardware (físico) & Khepera~IV nº~0471 (CETELI/UFAM) --- transferência futura. \\
Humano & Equipe de 4 integrantes (ver Tabela~2). \\
\bottomrule
\end{tabularx}
\end{table}

\subsection{Esforço estimado}
Estima-se esforço total da ordem de \textbf{120--160 homens-hora}, distribuído entre requisitos (15\%), projeto (20\%), implementação (35\%), verificação (20\%) e gerência/documentação (10\%).

% =====================================================================
\section{Gerência de riscos}
% =====================================================================
A Tabela~\ref{tab:riscos} apresenta os principais riscos, sua exposição (probabilidade $\times$ impacto) e estratégias de mitigação/contingência.

\begin{table}[H]\centering
\caption{Registro de riscos do projeto.}
\label{tab:riscos}
\renewcommand{\arraystretch}{1.2}
\begin{tabularx}{\textwidth}{@{}l X c c X@{}}
\toprule
\textbf{Id} & \textbf{Risco} & \textbf{Prob.} & \textbf{Imp.} & \textbf{Mitigação / Contingência} \\
\midrule
R1 & Aprisionamento em mínimo local (robô preso em obstáculo) & Média & Alto & FSM com contorno comprometido (Bug); contingência: planejamento global. \\
R2 & Descalibração de sensores (limiar errado) & Média & Alto & Calibração empírica via telemetria; contingência: reajuste de limiares. \\
R3 & Estouro do período do laço (perda de \emph{deadline} RT) & Baixa & Alto & WCET enxuto, sem alocação dinâmica no laço; contingência: aumentar período. \\
R4 & Divergência simulação $\to$ robô real (uso de GPS) & Alta & Médio & Isolar camada de localização; substituir por odometria; documentar limitação. \\
R5 & Atraso de cronograma/integração & Média & Médio & Entregas incrementais; folgas; replanejamento. \\
R6 & Indisponibilidade do robô físico & Média & Médio & Desenvolvimento e validação em simulação. \\
\bottomrule
\end{tabularx}
\end{table}

% =====================================================================
\section{Gerência de configuração e qualidade}
% =====================================================================
\textbf{Configuração:} controle de versão com Git; versionamento dos artefatos de documentação (campo ``Controle de versões''); separação clara entre mundo (\texttt{.wbt}), controlador (\texttt{.c}) e documentos. \textbf{Qualidade:} revisões por pares a cada entrega; validação por testes automatizáveis em modo \emph{headless}; rastreabilidade requisitos--testes (ver Análise de Requisitos e Plano de Testes).

% =====================================================================
\section{Monitoramento e controle}
% =====================================================================
O progresso é acompanhado pelos marcos da Tabela~4. Desvios de cronograma ou de qualidade são tratados em reuniões de revisão ao fim de cada entrega, com atualização deste plano (incremento de versão). A telemetria instrumentada do controlador funciona como métrica objetiva de progresso técnico.

\end{document}
